\documentclass[11pt]{amsart}
\usepackage[margin=1.25in]{geometry}
\usepackage{amsmath,amssymb,amsthm}
\usepackage{booktabs}
\usepackage[colorlinks=true,linkcolor=blue,citecolor=blue,urlcolor=blue]{hyperref}

\newtheorem{theorem}{Theorem}[section]
\newtheorem{proposition}[theorem]{Proposition}
\newtheorem{finding}[theorem]{Adverse Finding}
\theoremstyle{definition}
\newtheorem{definition}[theorem]{Definition}
\newtheorem{ruleenv}[theorem]{Rule}
\theoremstyle{remark}
\newtheorem{remark}[theorem]{Remark}

\begin{document}

\title[CAS PCBA 8: adaptive cognitive circuit engineering]{CAS PCBA 8: adaptive cognitive circuit engineering with typed, fail-closed design rules}
\author{Little Henderson}
\thanks{White paper v0.1, 2026-08-07. Artifact status: \textsc{draft / open / ring 4 / not licensed / held}. Engineering baseline: \texttt{CAS\_PCBA\_Engineer.pdf}, 743~pp., source ID \texttt{PCBA-ENG-20260726}, SHA-256 \texttt{294320B3\ldots}. Controlling artifacts: source audit and spec v0.1 (SHA-256 \texttt{19707B92\ldots}); CCDL construct coverage v0.1 (SHA-256 \texttt{5BB797EC\ldots}); publication delta Paper~6 at base commit \texttt{a1cd9901}. Implementation \texttt{src/cas\_pcba\_v8/}. Drafting this paper promotes no claim.}
\date{August 7, 2026}

\begin{abstract}
CAS PCBA 8 is a typed, provenance-bearing research architecture for describing and checking circuit-shaped cognitive-system models. ``Circuit'' names an engineering organization---typed components, ports, nets, state annotations, calibration references, deterministic traces---and does not assert a physical electrical implementation or identify model quantities with voltage, current, or resistance. The research question is whether a circuit-shaped modeling language can make cognitive-system hypotheses more explicit, type-safe, replayable, and falsifiable \emph{without} converting an analogy into a physical, causal, safety, or product claim. The candidate supports a qualified affirmative at the software-contract level: semantic identities remain typed; board structure is canonicalized; a deterministic design-rule checker rejects incompatible signal/unit pairs, unresolved ports, unsafe feedback declarations, missing calibrated probes, and unreachable fuse thresholds; supported simulations emit deterministic provenance-rich traces; and only certified finite, deterministic, stateless transitions may cross the TKS~v8 adapter boundary. Two adverse source findings are preserved rather than silently repaired: the published formal CCDL grammar does not cover the book's own five-board reference corpus (four direct syntactic counterexamples), and the published Cognitive Safety Board clamps a red stage at $0.1$ upstream of a fuse threshold of $0.9$, making the stated crossing unreachable under the declared clamp semantics. At the bound Gate~4 receipt, 82 CAS PCBA v8 test methods pass within an integrated 208-method three-suite run. These receipts establish tested software behavior only---not language completeness, empirical cognitive validity, physical materiality, safety effectiveness, or product readiness.
\end{abstract}

\maketitle

\section{What ``circuit'' means here}

CAS PCBA treats a candidate cognitive pathway as a directed, typed network whose elements must declare interfaces, allowed values, transfer behavior, statefulness, time semantics, calibration references, and evidence limits. Three separations are load-bearing.

\begin{enumerate}
\item \textbf{Model versus world.} A simulated transfer rule is behavior of the declared software profile, not evidence that a cognitive or physical system obeys it.
\item \textbf{Design value versus measurement.} A locally supplied parameter or calibration label remains structural metadata until authenticated measurement evidence and uncertainty are supplied.
\item \textbf{Analogy versus identity.} Borrowed circuit names do not make CAS quantities SI electrical quantities.
\end{enumerate}

The engineering literature supports this discipline at the level of method, not CAS validity: VHDL for machine-readable formal hardware description \cite{vhdl}, Statecharts for explicit state, hierarchy, and concurrency semantics \cite{harel}, and the verification-versus-validation distinction for simulation models \cite{sargent}.

\section{Typed semantic and board model}

\begin{definition}[State types]\label{def:states}
Building on the CAS Core 8 semantic-identity kernel, the candidate keeps distinct:
\begin{align*}
\mathrm{StateLiteral} &= \mathrm{Noetic}\times\mathrm{Domain}\times\mathrm{Drive}_{\mathrm{opt}}\times\mathrm{Register}_{\mathrm{opt}},\\
\mathrm{NDCell} &= \mathrm{Noetic}\times\mathrm{Domain},
\end{align*}
and $\mathrm{SimplexAddress}_{40}$, a strictly positive normalized profile over the forty canonical ND-cells.
\end{definition}

\begin{remark}
The projection $\mathrm{StateLiteral}\to\mathrm{NDCell}$ forgets Drive and Register and is explicitly many-to-one---not an isomorphism. A simplex address is neither a one-hot cell nor an element of the finite TKS kernel. The mixed-simplex geometry retains a named Fisher--Rao profile \cite{amari}; the disclosed binary64 implementation is a numerical pseudometric at its resolution, not a proof that all distinct profiles are separated. Projector geometry is a separate optional research profile that caps and versions its tolerance, records validation residuals, and binds matrix, basis, dimension, rank, and tolerance in one digest.
\end{remark}

\begin{ruleenv}[Values and arithmetic]
Finite carrier values use tagged exact type-plus-value identity, so integer $1$, Boolean \emph{true}, and string \texttt{"true"} are distinct throughout design-rule checking, execution, traces, and adapters. All decimal operations run under one immutable, versioned context created from primitive constants; ambient caller precision is irrelevant, and the context identity enters reports and results.
\end{ruleenv}

A board contains immutable, identity-bearing records for components and kinds; typed ports (direction, signal kind, unit); typed parameters and finite-state annotations; nets and endpoints; test points and calibration references; declared feedback loops; and source and research-boundary metadata. Board identity includes the complete canonical structure.

\section{The design-rule checker}

The checker computes one deterministic diagnostic report under a closed catalog. Errors fail the board; warnings never become an implicit clean result. Checks cover five classes.

\begin{enumerate}
\item \emph{Structural integrity}: unique nonblank identities; closed component kinds and exact parameter keys; missing or unknown semantics fail closed.
\item \emph{Typed connectivity}: every endpoint resolves; exactly one typed driver per net; driver and load signal kinds and units agree under a closed admissibility table; required inputs connect or have an executable default in the selected simulator profile.
\item \emph{Calibration and probe integrity}: executable test points require versioned calibration references with profile-digest shape, instrument and protocol identity, uncertainty, measurement kind, and net-unit compatibility; a locally constructed calibration record cannot make itself evidence-eligible; safety-marked paths require marked probes, but a probe's presence does not establish safety.
\item \emph{Feedback analysis}: a directed component graph with strongly connected components \cite{tarjan}; each realizable declared cycle needs an exact feedback declaration, a declared unit delay, an existing dimensionless calibration, and a conservative sub-unity gain bound; unsupported or resource-bounded analysis yields an unresolved-safety diagnostic, never a pass.
\item \emph{Threshold reachability}: conservative upper bounds propagate through causes, bias and gain stages, clipping, and cutoffs; a safety-critical fuse whose threshold exceeds the resolved upstream bound is flagged unreachable; an unresolved bound is not treated as reachable.
\end{enumerate}

\section{Adverse source findings}

\begin{finding}[CCDL grammar does not cover its own corpus]\label{find:grammar}
The consolidated EBNF admits only seven \texttt{circuit\_item} alternatives and only layer or via declarations in a board block, while the five published reference boards additionally use I/O, CJT, component, rail, dial, \texttt{use}, monitor, event, verification, and other forms listed in prose but absent from the grammar. The machine-readable coverage record documents four direct syntactic counterexamples and one bounded listing without an established counterexample. Pinpoint evidence: additional Part~IV statement leaders reserved at PDF page~578; missing appendix forms called errata at PDF pages 591--592; the consolidated EBNF at PDF pages 598--599; reference listings at PDF pages 606--632.
\end{finding}

\begin{ruleenv}
The v0.1 implementation declares only its programmatic Python model and does not claim to parse the complete CCDL book language. Any later grammar must publish an exact construct-coverage matrix, reject unsupported forms, and round-trip every admitted board canonically.
\end{ruleenv}

\begin{finding}[The CSB fuse cannot trip]\label{find:csb}
The published Cognitive Safety Board clamps its red stage at $\mathrm{sat}=0.1$ while the downstream fuse threshold is $0.9$; under the stated clamp semantics the fuse cannot trip. The same listing uses non-finite gain $\beta=\infty$ despite finite grammar and validator rules, and the book defines no fail-closed latch state machine showing that an open fuse preserves rather than suppresses the refusal path. Evidence: red-stage and fuse values at PDF page~342, repeated at PDF page~608; the Appendix~D note at PDF page~610 replaces the infinite gain with a large finite value but retains the clamp and does not reconcile it with the threshold.
\end{finding}

\begin{proof}[Justification under clamp semantics]
Saturation is a plateau: the red-stage output is bounded above by its clamp value $0.1$. Every downstream stage in the declared path is monotone and non-amplifying under the published parameters, so the fuse input never exceeds $0.1<0.9$. Hence the threshold crossing is unreachable.
\end{proof}

\begin{remark}
The unmodified parameters are retained as a known-fault fixture that the design-rule checker must flag. Any repair is a separately versioned design requiring threshold reachability, finite parameters, latch state transitions, fuse truth tables, deterministic traces, hazard analysis, and independent safety review. This is valuable negative evidence about the source and a regression guard for the checker---not a validation of the board's safety function.
\end{remark}

\begin{remark}[Claim quarantine]
The following remain hypotheses or prohibited product inferences: exact thought/transistor isomorphism; shared four-region identity; overload-induced semantic involution; Fisher--Rao causal gain; physical electrical identity for CAS quantities; semantic correctness from a response-region label; safety effectiveness, certification, regulatory or clinical use; complete compiler conformance; production readiness; and disclosure or license authority from a self-hash or generated receipt.
\end{remark}

\section{Narrow deterministic simulation}

Simulation is admitted only for supported component profiles with declared ports, transfer rules, clipping, statefulness, and time semantics. The current profile topologically orders the non-delay graph, uses declared unit-delay behavior to break admitted feedback, and emits a canonical trace binding board and component identities, the exact design-rule report, component catalog and decimal-context identities, simulator profile and source baseline, typed inputs and outputs, component events and tick linkage, test-point observations, measurement IDs with parent/child derivations, and the fixed ring-4 boundary. Measurement construction rechecks child existence, derivation acyclicity, tick linkage, units, measurement kinds, and digest form---a discipline resembling the entity/activity/derivation separation of W3C PROV \cite{prov} without claiming PROV-O conformance.

The correct reading of a green fixture is: \emph{the software performed the declared typed transition and produced the declared trace for the exact fixture under the exact profile}---never that the fixture proves a cognitive mechanism, causal relation, physical law, safe refusal, or product capability.

\section{TKS v8 and Noetic Topos 2 boundaries}

Only a certified total, finite, deterministic, stateless, non-stochastic component transition may be adapted to a TKS v8 finite arrow. Chains become semigroupoid paths; only an endomap has powers. The adapter certificate binds the exact TKS kernel bytes, board, component, transition table, tagged carrier, semantic profile, and research boundary; ambient module substitution is rejected. The adapter proves at most type-correct representation in the finite reference kernel---it carries no automatic CAS semantic theorem. Continuous simplex profiles, response curves, instrumented observations, stateful devices, stochastic kernels, and memristive history variables remain outside the finite arrow and require distinct types and comparison maps.

From Noetic Topos 2.0 the candidate inherits the ambient-versus-resolved evidence distinction: raw or partially resolved records may be treated as external evidence associated with the ambient construction, but a marked cover does not make a measurement true, and records do not enter resolved semantics until restriction and gluing under one calibration profile are defined and tested.

\section{Bounded software evidence}

At the bound commit, the integrated Gate~4 receipt records 54 hash-bound artifacts and three green focused suites:

\begin{center}
\small
\begin{tabular}{@{}lrl@{}}
\toprule
Suite & Methods & Test-ID SHA-256 (prefix) \\
\midrule
CAS Core v8 & 61 & \texttt{E2AF5014} \\
CAS PCBA v8 & 82 & \texttt{D5F776C5} \\
CAS-Ohm v1 & 65 & \texttt{2E81A685} \\
\bottomrule
\end{tabular}
\end{center}

The combined run comprises 208 test methods and 272 parameterized subtests, with 11 claim records all ring~4 and not licensed, 53 review criteria (20 open), four direct EBNF counterexamples, and eight CAS-Ohm source conflicts. The repository-wide collector was not clean at the bound receipt (four pre-existing package-path namespace collisions after 1{,}200 collected items), so the focused suites cannot support a repository-wide verification claim.

\section{Threats to validity}

\emph{Construct}: no experiment establishes that the selected CAS types correspond to cognitive mechanisms or latent model structure. \emph{Internal}: bounded regression suites, not proofs of all paths; property-based and fuzz campaigns \cite{quickcheck}, independent implementations, and a separately sourced mutation-testing campaign remain open; the full CCDL parser does not exist. \emph{External}: no live calibrated model, physical board, instrument chain, held-out replication, or multi-model evaluation is in the admitted evidence. \emph{Numerical}: independent audit of the decimal profile and the binary64 Fisher--Rao pseudometric remains open. \emph{Causal}: connectivity, gain composition, prediction, and residuals do not establish causation; causal claims require explicit models, identification assumptions, and interventions \cite{pearl}; none is admitted. \emph{Safety and governance}: Finding~\ref{find:csb} blocks any safety-effectiveness inference; hazard analysis, independent safety engineering, regulatory and clinical review, authenticated claim authority, canon admission, disclosure approval, and product acceptance are all open.

\section{Dependency-ordered research program}

(1) Specify a complete versioned CCDL grammar and a golden five-board corpus. (2) Implement canonical parse/serialize round trips; fail closed on unsupported syntax. (3) Define separate capacitor, PLL, memory, power, feedback, stochastic, and memristive state and time profiles. (4) Run a property-based and random-generation campaign over checker and simulation invariants \cite{quickcheck}; separately specify any mutation-testing campaign. (5) Reproduce the exact source and test bundle in a hermetic toolchain. (6) Preregister calibration and held-out validation through the Gate~5 control kernel without promoting Gate~5 status. (7) Require executor/auditor numerical agreement, falsifiers, negative outcomes, and independent replication. (8) Complete internal cross-review against stable hashes; then commission external formal, systems, experimental, safety, and claim-language review.

\section{Conclusion}

CAS PCBA 8's value is not that it proves cognition is electrical---it does not and cannot at this stage. Its value is that it turns informal pathway diagrams into typed, inspectable, versioned software objects and \emph{preserves negative evidence} when a source or design cannot satisfy its own contract: a grammar that does not cover its own corpus, and a safety board whose fuse cannot trip, are recorded as first-class findings with page-anchored evidence rather than patched away. The strongest current result is architectural: every analogy is forced through a typed interface, and every executable claim is bound to a profile, source, calibration record, and trace.

\begin{thebibliography}{9}
\bibitem{vhdl} IEEE Std 1076-2019, \emph{VHDL Language Reference Manual}.
\bibitem{harel} D.~Harel, ``Statecharts: a visual formalism for complex systems,'' \emph{Sci.\ Comput.\ Program.}\ 8 (1987), 231--274.
\bibitem{tarjan} R.~Tarjan, ``Depth-first search and linear graph algorithms,'' \emph{SIAM J.\ Comput.}\ 1 (1972), 146--160.
\bibitem{prov} W3C, \emph{PROV-O: The PROV Ontology}, 2013.
\bibitem{sargent} R.~G. Sargent, ``Verification and validation of simulation models,'' \emph{J.\ Simulation} 7 (2013), 12--24.
\bibitem{quickcheck} K.~Claessen and J.~Hughes, ``QuickCheck: a lightweight tool for random testing of Haskell programs,'' \emph{ICFP} (2000).
\bibitem{amari} S.-i.~Amari, \emph{Information Geometry and Its Applications}, Springer, 2016.
\bibitem{pearl} J.~Pearl, \emph{Causality}, 2nd ed., Cambridge University Press, 2009.
\end{thebibliography}

\end{document}
